\documentclass[12pt, a4paper]{article}
\usepackage[spanish]{babel}
\usepackage{color,graphicx}
\usepackage{graphics}
\usepackage{amsmath,amssymb,amsfonts,amsthm}
\usepackage[T1]{fontenc}
\usepackage{float}
\usepackage{xcolor}
\usepackage{amsthm}
\usepackage{tikz}
\usepackage{enumerate}
\usepackage{picinpar}
\usepackage{newtxtext,fourier}
\makeatother
\usepackage{thmtools}
\usepackage[framemethod=TikZ]{mdframed}
\mdfsetup{skipabove=1em,skipbelow=1em}

\newcommand{\sqed}{\hfill\includegraphics[height=2ex]{sancoya.png}}

\theoremstyle{definition}
\declaretheoremstyle[
    headfont=\bfseries\color{black!70!black}, bodyfont=\normalfont,
    mdframed={
        linewidth=0.7pt,
        rightline=true, topline=true, bottomline=true,
        linecolor=black, backgroundcolor=black!2!white,
    }
]{thmbox}

\declaretheorem[style=thmbox, name=Definición]{definition}
\declaretheorem[style=thmbox, numbered=no, name=Ejemplo]{example}
\declaretheorem[style=thmbox, name=Proposición]{prop}
\declaretheorem[style=thmbox, name=Teorema, numbered=no]{theorem}
\declaretheorem[style=thmbox, name=Lema]{lemma}
\declaretheorem[style=thmbox, numbered=no, name=Corolario]{corollary}
\declaretheorem[style=thmbox, numbered=no, name=Nota]{note}




\usepackage{geometry}
 \geometry{
 a4paper,
 total={170mm,260mm},
 left=25mm,
 top=30mm,
 }
\begin{document}
\setlength{\parindent}{0cm}
\hoffset-0.46cm
\voffset-1.46cm
\noindent
\hbox to \linewidth{
  \begin{minipage}[t]{40pt}
    \vspace*{0pt}
\includegraphics[scale=0.22]{Graphics/logo.eps}
  \end{minipage}
  \hss
  \begin{minipage}[t]{0.8\linewidth}
    \vspace*{15pt}\centering
     Universidad Nacional de Colombia\\
    \small Facultad de Ciencias\\
    Teoría analítica de números\\
    Tarea II
  \end{minipage}
  \hss
}
\phantom{xd}\\
\small
\textsf{Mateo Andrés Manosalva Amaris}
\normalsize
\dotfill
\vspace{0.3cm}

\textbf{Ejercicio 1.}
\begin{theorem}
    Dado $x\geq 2$,
    $$\psi(x)-\vartheta(x)=O(\sqrt{x}\log x)$$
\end{theorem}

\begin{proof}
Por definición 
\begin{align*}
\psi(x)-\vartheta(x)&=\sum_{p^k\leq x}\log p-\sum_{p\leq x}\log p\\
&=\sum_{\substack{p^k\leq x\\k\geq 2}}\log p   
.\end{align*}

Note que, en realidad, la suma es finita, dado $p$ primo, siempre existe $k$ tal que $p^k>x$. Más precisamente, dado $p$, la suma para cuando $k=\log_p(x)$, obtenemos lo siguiente:

$$
\psi(x)-\vartheta(x)=\sum_{p^k\leq x}\log p \sum_{2\leq k\leq\log_p(x)}1
,$$
podemos escribir esto de una manera más sencilla de controlar, en efecto, si $p$ es un número primo tal que $x^{\frac{1}{k}}<p<x^{\frac{1}{2}}$, entonces $p^k>x$, por lo que realmente, podemos tomar la primera suma hasta $p\leq \sqrt{x}$ y esto no altera la igualdad. Esto es

\begin{align*}
  \psi(x)-\vartheta(x)&=\sum_{p\leq \sqrt{x} }\log p \sum_{2\leq k\leq\log_p(x)}1\\
  &\leq\sum_{p\leq \sqrt{x} }\log p \left[\dfrac{\log x}{\log p} \right]\\
  &\leq \sqrt{x}\log x
.\end{align*}
\end{proof}

\textbf{Ejercicio 2.}
\begin{theorem}
    Dado $x\geq 2$,
    $$\psi(x)-\vartheta(x)=O(\sqrt{x})$$
\end{theorem}

\begin{proof}
  En efecto

  \begin{align*}
    \psi(x)-\vartheta(x)&=\sum_{2\leq k\leq \log_p(x)}\vartheta(x^{1/k})\\
    &\leq\vartheta(x^{1/2})+\vartheta(x^{1/3})\sum_{3\leq k\leq \log_p(x)}1\\
    &\leq\vartheta(x^{1/2})+\vartheta(x^{1/3})\left[\frac{\log x}{\log p}\right]
  .\end{align*}
Por los teoremas de Chevyshev, sabemos que $\vartheta(x)=O(x)$, así
$$
\psi(x)-\vartheta(x)=O(\sqrt{x})+O(x^{1/3}\log x)=O(\sqrt{x})
$$
ya que
$$\lim_{x \to \infty} \dfrac{x^{1/3}\log x}{\sqrt{x}}=\lim_{x \to \infty}x^{-1/6}\log x=0$$
\end{proof}

\end{document}